% Transcription — fonds Alexandre Grothendieck, shelfmark 11, batch 1 (pages 1-20).
% Established from the facsimile archives/batches/11/batch-01.pdf (a mirror of
% https://grothendieck.umontpellier.fr/11.pdf), per the skill
% transcribe-grothendieck. The batch file carries no cover sheet: PDF page N is
% archivists' page N (the pencilled numbers confirm it on every leaf).
%
% Pass: Opus 5 (claude-opus-5), 2026-09-14 — first pass, unchecked against the pages by a human.
%
% Revised 2026-09-23 (Opus 5.5), after the find-novelty pass. Header only: the
% typescript headings were listed as twelve (-2.24- to -2.33-, -3.1-, -3.2-)
% for nine typescript pages; re-read on the facsimile, the nine pages carry
% nine headings, given below. The transcription itself is unchanged.
%
% What the batch is. Ten leaves in his hand (the even pages 2-20), written on
% the backs of a typescript in English on geometric invariant theory: the nine
% odd pages 3-19 are that typescript, upside down and without a mark of his,
% and page 1 is blank. The typescript's page numbers run backwards through the
% batch: -3.2- (page 3), -3.1- (5), -2.33- down to -2.28- (7-17), -2.26- (19);
% there is no -2.27-. They are skipped without further mention. The leaves carry no date and no
% pagination of his own.
%
% The mathematics is one continuous run: the formalism of correspondences
% attached to a bivariant theory. V is a category with products of two
% objects; A : X ↦ A(X) a functor to a category C of (graded) « bivariant »
% algebras, i.e. with f_* and f^*, subject to an exchange (base change)
% condition and the projection formula. He builds a new category V^A with the
% objects of V and Hom_{V^A}(X,Y) = A(X×Y), composition
% v∘u = p_{31*}(p_{32}^*(v) p_{21}^*(u)), and two functors φ : V → V^A and
% ψ : V^A → k-mod gr. with ψφ = ωA.
%
%   2   the data; V^A, φ, ψ and the square they make commute; Ob V^A = Ob V;
%       the composition law, boxed.
%   4   « Associativité »: w(vu) and (wv)u both computed down to
%       p_{XT*}(p_{ZT}^*(w) p_{YZ}^*(v) p_{XY}^*(u)), by exchange and
%       projection, with two diagrams of projections from X×Y×Z×T.
%   6   « Unités »: Δ_X = A(δ_X)_*(1_{A(X)}); u∘Δ_X = u through the cartesian
%       square of δ_X × id_Y; a boxed marginal on the signs (−1)^{ux}; the NB
%       that composition needs only projections, units the diagonals too.
%   8   the two cases of exchange/projection actually used; definition of
%       φ_{X,Y}(f) = A(Γ_f)_*(1) through the graph Γ_f = (id_X, f).
%  10   φ is a functor: φ(id_X) = Δ_X, and φ(g)φ(f) = φ(gf) through a
%       cartesian square of graphs and ν(x) = (x, f(x), gf(x)); then ψ:
%       u(x) = p_{2*}(u p_1^*(x)).
%  12   ψ is a functor: Δ_X(x) = x, and v(u(x)) = (v∘u)(x), both sides
%       reduced to p_{Z*}(p_{YZ}^*(v) p_{XY}^*(u) p_X^*(x)).
%  14   commutativity of the square, f_*(x) = φ(f)(x); « Remarque » on
%       interpreting the contravariant side A(f)^*, answered by φ(f)^* obtained
%       by transposition; the canonical s on V^A, u ↦ ^t u.
%  16   φ* = sφ°, ψφ* = ω*A°; the case of a final object e with A(e) = k:
%       the A(X) « jaugés » (read doubtfully) by the forms k_X; if they are
%       Poincaré algebras, u_* and u^* are transposes, computed « mod. signe ».
%  18   « Remarque »: only the monoid structure of the A(X) was used, so the
%       theory should be written for « monoïdes bivariants »; a cancelled
%       passage on signs; the proposal of graded monoids (A^i)_{i∈Z} with
%       involutions σ^i and a degree d(A^*) fixing the degree of f_*.
%  20   sub-objects A'(X) ⊂ A(X): equivalence of 1° stability under inverse
%       and direct images and products, and 2° a) u(x) ∈ A', b) x ⊠ y ∈ A',
%       c) graphs in A'; a struck and rewritten NB (the A' as subcategories
%       of V^A); the proof both ways, with boxed marginal tags.
%
% Notation fixed for the batch. His curled V is \mathcal{V}, and the new
% category \mathcal{V}^{A}; the functor A is sometimes drawn script, sometimes
% roman, and is given as roman A throughout; C, e, k stay roman. Pull-backs and
% push-forwards as p^{*}, p_{*}, with the indices he writes (p_{31}, p_{XZ},
% q_1, r_2 …) left as they stand even where they disagree with a diagram's
% labels (page 2). The external product is \boxtimes, the internal \otimes.
%
% Material facts stated once. The hand is fast throughout and the connective
% prose often does not carry while the formulas do; page 18 is cancelled in a
% block (a struck line, a boxed passage and a long wavy stroke), and page 20
% has a struck NB rewritten in a box. Duplicate draft sketches of a diagram
% (page 10, top left) are noted rather than redrawn.
%
% Batch boundary. The proof on page 20 closes on the leaf (« Stabilité par ⊠
% claire. ») and the lower half is blank, so no sentence crosses into batch 2;
% but the subject of the A' is taken up again on page 22 (batch 2), and a \note
% at the end says so.

\documentclass[11pt,a4paper]{article}
% Le préambule commun à toutes les transcriptions du fonds Grothendieck.
%
% Il tient en une idée : **l'incertitude se compose, elle ne se résout pas.**
% Une transcription qui lisse un mot douteux a détruit la seule chose pour
% laquelle on la faisait — un lecteur doit pouvoir savoir, à chaque endroit,
% ce qui était sur le feuillet et ce qui a été ajouté. Les six macros qui
% suivent sont donc l'appareil critique, et non de la décoration.
%
% Les mêmes commandes sont reconnues par `scripts/render.mjs`, qui produit la
% vue de lecture du site. Une macro ajoutée ici doit l'être aussi là-bas,
% faute de quoi le rendu échouera bruyamment — ce qui est voulu : un rendu
% silencieusement faux est pire qu'un rendu absent.

\ProvidesPackage{grothendieck}[2026/08/08 Transcriptions du fonds Grothendieck]

\usepackage{amsmath,amssymb,amsthm}
\usepackage{mathrsfs}
\usepackage{stmaryrd}
% Les diagrammes commutatifs sont partout dans ce fonds : c'est la langue
% même de la géométrie algébrique de Grothendieck, et une transcription qui
% les rendrait en prose ne transcrirait rien.
\usepackage{tikz-cd}
% Français par défaut — c'est la langue des feuillets — mais l'anglais est
% chargé aussi : la traduction et le résumé sont des documents anglais, et la
% typographie française (espaces avant les deux-points) y serait une faute.
% Ces documents basculent par \selectlanguage{english} en tête de corps.
\usepackage[english,french]{babel}
\usepackage{fontspec}
\usepackage{geometry}
\geometry{a4paper,margin=25mm,marginparwidth=28mm}
\usepackage{marginnote}
\usepackage{xcolor}
% ulem, not soul. `soul`'s \st and \ul are built on a character-by-character
% scan that XeLaTeX's font machinery breaks: the document fails with an
% undefined control sequence rather than degrading. ulem does the same two jobs
% and is Unicode-engine safe. `normalem` stops it hijacking \emph, which the
% transcriptions use for Grothendieck's own emphasis.
\usepackage[normalem]{ulem}
\usepackage{hyperref}
\hypersetup{colorlinks=true,linkcolor=black,urlcolor=black,pdfborder={0 0 0}}

% --- Métadonnées du lot -----------------------------------------------------
% Reprises telles quelles par le script de rendu, qui les affiche en tête de
% la vue de lecture. Elles fixent ce que le fichier prétend couvrir : sans
% elles, une transcription flotte sans point d'attache dans un fonds de seize
% mille feuillets.

\newcommand{\folder}[1]{\gdef\grothFolder{#1}}
\newcommand{\batch}[1]{\gdef\grothBatch{#1}}
\newcommand{\leaves}[2]{\gdef\grothFirst{#1}\gdef\grothLast{#2}}
\newcommand{\foldertitle}[1]{\gdef\grothFoldertitle{#1}}
\gdef\grothFolder{?}\gdef\grothBatch{?}\gdef\grothFirst{?}\gdef\grothLast{?}\gdef\grothFoldertitle{}

% --- Le résumé --------------------------------------------------------------
%
% Il ouvre la lecture modernisée, sous le titre « Résumé », et il est composé
% en petit. Ce n'est pas un ornement : c'est le seul passage du document qui
% ne soit pas des mathématiques, et un lecteur doit pouvoir le reconnaître
% avant de l'avoir lu — puis le sauter s'il connaît déjà le sujet, ou s'y
% arrêter s'il ne le connaît pas. Le filet qui le suit marque la reprise du
% corps.

\newenvironment{resume}
  {\small\setlength{\parskip}{0.45em}}
  {\par\medskip\hrule\medskip}

% --- Mots-clés ---------------------------------------------------------------
%
% En anglais, en clôture du résumé de la lecture modernisée : le vocabulaire
% moderne sous lequel chercher ce que les feuillets construisent. Le manifeste
% du site (`npm run manifest`) les extrait pour étiqueter la cote entière —
% cette ligne est la source unique des tags, il n'y a pas de fichier à côté.
\newcommand{\keywords}[1]{\par\smallskip\noindent
  {\footnotesize\textsc{Keywords} --- \textit{#1}}\par}

% --- L'appareil critique ----------------------------------------------------

% Le numéro de feuillet, en marge. C'est l'ancre qui permet de régler un
% désaccord sur une formule en regardant *un* feuillet plutôt qu'une chemise
% de six cents. Le site s'en sert aussi pour tourner les pages du fac-similé
% quand on fait défiler la transcription.
\newcommand{\leaf}[1]{%
  \marginnote{\footnotesize\color{gray!70}\textsf{#1}}%
  \label{leaf:#1}\ignorespaces}

% Illisible : on ne devine pas. Un mot inventé qui se lit comme les autres est
% le pire résultat possible.
\newcommand{\ill}{\textcolor{red!60!black}{[\,\ldots\,]}}

% Lecture douteuse : le mot est donné, et signalé comme incertain.
\newcommand{\uncertain}[1]{\uline{#1}}

% Ajout de l'éditeur — une lettre manquante, un mot sous-entendu.
\newcommand{\add}[1]{\textcolor{blue!50!black}{[#1]}}

% Biffé par Grothendieck lui-même. Ce qu'il a rayé fait partie du document :
% une rature dit souvent où la pensée a changé de direction.
\newcommand{\struck}[1]{\sout{#1}}

% Note du transcripteur, sur l'état matériel du feuillet.
\newcommand{\note}[1]{\par\smallskip{\small\color{gray!60!black}\textit{[#1]}}\par\smallskip}

% Note marginale de Grothendieck, distincte du corps du texte.
\newcommand{\marginal}[1]{\par\smallskip\noindent\hspace{1em}%
  {\small\color{gray!60!black}#1}\par\smallskip}

% --- Titre ------------------------------------------------------------------

\AtBeginDocument{%
  \begin{center}
    {\large\bfseries Fonds Alexandre Grothendieck --- cote n\textsuperscript{o}~\grothFolder}\\[2pt]
    {\itshape\grothFoldertitle}\\[4pt]
    {\small feuillets \grothFirst--\grothLast\ (lot \grothBatch)}\\[2pt]
    {\footnotesize\color{gray!60!black}Transcription établie d'après le fac-similé numérisé
     par l'Université de Montpellier.\\
     Les lectures incertaines sont soulignées, les illisibles notées [\,\ldots\,],
     les ajouts entre crochets.}
  \end{center}
  \vspace{1em}}

\endinput


\folder{11}
\batch{1}
\pages{1}{20}
\dating{[à partir de 1961]}
\watermark{Édition de démonstration}
\foldertitle{Formalisme des correspondances : notes manuscrites (s.d.)}

\begin{document}

\page{2}

$\mathcal{V}$ catégorie (objets $X, Y, \ldots$) avec produits de deux objets
(\uncertain{et un} objet final)

$A \colon X \mapsto A(X)$ foncteur cov.\ de $\mathcal{V}$ dans la catégorie $C$
des \uncertain{algèbres gr.}\ bivariantes, satisfaisant une condition d'échange
\struck{\ill{}} \uncertain{(au sens gr.)} (dont on \uncertain{aura} besoin
\struck{\ill{}} \ill{} \ill{} la suite ....).

\note{Le début de la ligne est surchargé : « gr. » est ajouté au-dessus, et un
mot biffé sous « algèbres » semble être « anticommutatives » ; la lecture de
cette ligne reste douteuse.}

On va définir une nouvelle catégorie, notée $\mathcal{V}^{A}$, et deux
foncteurs $\varphi, \psi$

\begin{tikzcd}
\mathcal{V} \arrow[r, "\varphi"] \arrow[dr, "A"'] & \mathcal{V}^{A} \arrow[r, "\psi"] & k\text{-mod gr.} \\
 & C \arrow[ur, "\omega"'] &
\end{tikzcd}

($\omega$ : oubli de la structure d'algèbre), rendant le carré ci-dessus
commutatif. De plus, le foncteur $\varphi$ induira une bijection (en fait, une
identité) sur les ensembles d'objets \struck{\ill{}} :
\[ \operatorname{Ob} \mathcal{V}^{A} = \operatorname{Ob} \mathcal{V} . \]

On pose
\[ \operatorname{Hom}_{\mathcal{V}^{A}}(X,Y) = A(X \times Y) . \]
On définit le produit de composition ainsi :
\[
\begin{array}{ccc}
\operatorname{Hom}_{\mathcal{V}^{A}}(X,Y) \times \operatorname{Hom}_{\mathcal{V}^{A}}(Y,Z) & \longrightarrow & \operatorname{Hom}_{\mathcal{V}^{A}}(X,Z) \\
(u, v) & \longmapsto & v \circ u
\end{array}
\]
\[ A(X \times Y) \times A(Y \times Z) \longrightarrow A(X \times Z) \]

\begin{tikzcd}
 & X \times Y \times Z \arrow[dl, "p_{12}"'] \arrow[dr, "p_{23}"] \arrow[dd, "p_{13}"] & \\
X \times Y & & Y \times Z \\
 & X \times Z &
\end{tikzcd}

\[ \boxed{\, v \circ u = p_{31*}\bigl( p_{32}^{*}(v)\, p_{21}^{*}(u) \bigr) . \,} \]

\note{Le diagramme porte les indices $p_{12}, p_{13}, p_{23}$, la formule
encadrée $p_{31}, p_{32}, p_{21}$ ; on les laisse tels quels.}

\page{4}

\subsection*{Associativité}

\[ X \xrightarrow{\;u\;} Y \xrightarrow{\;v\;} Z \xrightarrow{\;w\;} T \]

À prouver
\[ w(vu) = (wv)u \]
\struck{Le premier membre est égal à : \ill{}}

\begin{tikzcd}[column sep=small, row sep=small, nodes={font=\scriptsize}]
 & & X \times Y \times Z \times T \arrow[dl, "\alpha'"'] \arrow[drrr, "\nu'"] & & & & \\
 & X \times Y \times Z \arrow[dl, "\lambda"'] \arrow[d, "\mu"] \arrow[dr, "\nu"] & & & & X \times Z \times T \arrow[dlll, "\alpha"'] \arrow[d, "\gamma"] \arrow[dr, "\beta"] & \\
X \times Y & Y \times Z & X \times Z & & & X \times T & Z \times T
\end{tikzcd}

\note{Le but de $\mu$ est écrit $Y \times Z$ par-dessus une première lettre ;
une petite inscription le long de $\nu$ est \ill{}.}

\[
w(vu) \stackrel{\text{déf}}{=} \gamma_{*}\bigl( \beta^{*}(w)\, \alpha^{*}(vu) \bigr)
\stackrel{\text{déf de } vu}{=} \gamma_{*}\Bigl( \beta^{*}(w)\, \alpha^{*}\bigl( \nu_{*}( \mu^{*}(v)\, \lambda^{*}(u) ) \bigr) \Bigr)
\]
car $\alpha^{*} \nu_{*} = \nu'_{*} \alpha'^{*}$ [échange]
\[
= \gamma_{*}\Bigl( \beta^{*}(w)\, \nu'_{*}\bigl( \alpha'^{*}( \mu^{*}(v)\, \lambda^{*}(u) ) \bigr) \Bigr)
\]
avec
\[
\alpha'^{*}\mu^{*}(v)\, \alpha'^{*}\lambda^{*}(u) = p_{YZ}^{*}(v)\, p_{X,Y}^{*}(u)
\]
$\|$ formule de proj.\ (att.\ signes !)
\[
\nu'_{*}\bigl( \nu'^{*}\beta^{*}(w)\, p_{YZ}^{*}(v)\, p_{XY}^{*}(u) \bigr),
\qquad \nu'^{*}\beta^{*}(w) = p_{ZT}^{*}(w)
\]
\[
= p_{XT*}\bigl( p_{ZT}^{*}(w)\, p_{YZ}^{*}(v)\, p_{XY}^{*}(u) \bigr)
\]

\note{À gauche de ces lignes, un fragment biffé : \struck{$(-1)$ \ill{}}.}

\begin{tikzcd}[column sep=small, row sep=small, nodes={font=\scriptsize}]
 & & X \times Y \times Z \times T \arrow[dl, "n'"'] \arrow[dr, "b'"] & & \\
 & X \times Y \times T \arrow[dl, "a"'] \arrow[d, "c"] \arrow[dr, "b"] & & Y \times Z \times T \arrow[dl, "n"'] \arrow[d, "l"] \arrow[dr, "m"] & \\
X \times Y & X \times T & Y \times T & Y \times Z & Z \times T
\end{tikzcd}

\[
(wv)u = c_{*}\bigl( b^{*}(wv)\, a^{*}(u) \bigr)
= c_{*}\Bigl( b^{*}\bigl( n_{*}( m^{*}(w)\, l^{*}(v) ) \bigr)\, a^{*}(u) \Bigr)
\]
\note{Dans la seconde expression, un premier $m_{*}$ est biffé et remplacé par
$n_{*}$.}
formule d'échange :
\[
n'_{*}\bigl( b'^{*}( m^{*}(w)\, l^{*}(v) ) \bigr),
\qquad b'^{*}\bigl( m^{*}(w)\, l^{*}(v) \bigr) = p_{ZT}^{*}(w)\, p_{YZ}^{*}(v)
\]
\[
= c_{*}\Bigl( n'_{*}\bigl( p_{ZT}^{*}(w)\, p_{YZ}^{*}(v) \bigr)\, a^{*}(u) \Bigr)
\]
--- formule de proj.\ (att.\ signes !)
\struck{$= c_{*}$}
\[
= n'_{*}\bigl( p_{ZT}^{*}(w)\, p_{YZ}^{*}(v)\; n'^{*} a^{*}(u) \bigr),
\qquad n'^{*}a^{*}(u) = p_{XY}^{*}(u)
\]
\[
= p_{XT*}\bigl( p_{ZT}^{*}(w)\, p_{YZ}^{*}(v)\, p_{XY}^{*}(u) \bigr)
\]

\note{Devant $a^{*}(u)$ dans l'avant-dernière
ligne, un symbole est biffé et surchargé, lu $n'^{*}$ d'après l'accolade qui le
rend $p_{XY}^{*}(u)$. La dernière égalité est écrite dans la marge droite et
sort de la feuille.}

\page{6}

\subsection*{Unités}

On définit
\[ \Delta_{X} \in \operatorname{Hom}_{\mathcal{V}^{A}}(X,X) = A(X \times X) \]
en considérant
\[ \delta_{X} \colon X \to X \times X \qquad (\text{diag.}) \]
et définissant
\[ \Delta_{X} = A(\delta_{X})_{*}(1_{A(X)}) \]
\note{Après le signe $=$, un premier état est biffé : \struck{\ill{}}.}

Prouvons que $\Delta_{X}$ est une unité \uncertain{bilatère} pour la
composition : \uncertain{i.e.}
\[ X \xrightarrow{\;\Delta_{X}\;} X \xrightarrow{\;u\;} Y \]
\[ u \circ \Delta_{X} = p_{31*}\bigl( p_{32}^{*}(u)\, p_{21}^{*}(\Delta_{X}) \bigr) \]

\begin{tikzcd}
 & X \times X \times Y \arrow[dl, "p_{21}"'] \arrow[dr, "p_{32}"] \arrow[dd, no head] & \\
X \times X & & X \times Y \\
 & X \times Y &
\end{tikzcd}

\begin{tikzcd}
X \times Y \arrow[r, "\delta_{X} \times \mathrm{id}_{Y} = q"] \arrow[d, "p"'] & X \times X \times Y \arrow[d, "p_{21}"] \\
X \arrow[r, "\delta_{X}"] & X \times X
\end{tikzcd}

\[
p_{21}^{*}(\Delta_{X}) = p_{21}^{*}\bigl( \delta_{X*}(1) \bigr) \stackrel{\text{échange}}{=} q_{*}\bigl( p^{*}(1) \bigr) = q_{*}(1)
\]
donc
\[
p_{32}^{*}(u)\, p_{21}^{*}(\Delta_{X}) = p_{32}^{*}(u)\, q_{*}(1) = q_{*}\bigl( q^{*}( p_{32}^{*}(u) ) \bigr)
\]
(formule de proj., attention signe !)
\[ = q_{*}(u) \quad \text{car } p_{32}\, q = \mathrm{id}_{X \times Y} \]
donc
\[ u \circ \Delta_{X} = p_{31*}\, q_{*}(u) = u \]
car \struck{$q$} $p_{31}\, q = \mathrm{id}$.

\marginal{signes $(-1)^{ux}$ ; pour ces signes \uncertain{voir calculs} pour
$\Delta_{X} \circ v = v$, $v \colon Y \to X$}

On prouve de façon analogue que $\Delta_{X} \circ u = u$.

\emph{Donc on trouve bien une catégorie $\mathcal{V}^{A}$.}

[\emph{NB} Pour définir la composition on n'a eu besoin que des $p_{*}, p^{*}$
pour projections $X \times Y \rightrightarrows X$ ; mais pour
\struck{définir} construire les unités on a utilisé aussi $q_{*}, q^{*}$ pour
les imm.\ diagonales et leurs produits par un $Y$ ....]

\page{8}

De \uncertain{fait}, pour la \struck{définit} vérification de l'associativité,
on n'a utilisé la formule de projection et la formule d'échange que dans le cas

\begin{tikzcd}
 & X \times Y \times P \arrow[dl] \arrow[dr] & \\
X \times P \arrow[dr, no head] & & Y \times P \arrow[dl, no head] \\
 & P &
\end{tikzcd}

\note{Le $Y$ de $Y \times P$ est écrit par-dessus un $X$.}

mais pour vérifier que les \uncertain{$\Delta_{X}$} sont des unités, on a eu
besoin aussi du cas

\begin{tikzcd}
 & X \times Y \arrow[dl] \arrow[dr, hook] & \\
X \arrow[dr, hook, "\mathrm{diag.}"'] & & X \times X \times Y \arrow[dl, "\mathrm{proj}"] \\
 & X \times X &
\end{tikzcd}

Définissons le foncteur $\varphi$, en définissant pour tout
$X, Y \in \operatorname{Ob} \mathcal{V}$ une application
\[
\varphi_{X,Y} \colon \operatorname{Hom}_{\mathcal{V}}(X,Y) \longrightarrow \operatorname{Hom}_{\mathcal{V}^{A}}(X,Y) = A(X \times Y)
\]
\note{L'indice de $\varphi$ ressemble ici et plus bas à « $X,Z$ » ; le contexte
impose $X,Y$.}
Soit donc $f \colon X \to Y$ dans $\mathcal{V}$, on en déduit \ill{} graphe
\[ \Gamma_{f} = (\mathrm{id}_{X}, f) \colon X \to X \times Y \]
d'où
\[ A(\Gamma_{f})_{*} \colon A(X) \to A(X \times Y) \]
Posons alors \struck{\ill{}}
\[ \varphi_{X,Y}(f) = A(\Gamma_{f})_{*}(1) \]

\page{10}

Il faut prouver que $\varphi$ est bien un foncteur.

a) $\varphi_{X,X}(\mathrm{id}_{X}) = \Delta_{X}$. En effet,
\uncertain{c'est la définition} de $\Delta_{X}$ !

\note{Sur cette ligne, un premier état est biffé et la leçon « c'est la
définition » est écrite au-dessus.}

b) $X \xrightarrow{\;f\;} Y \xrightarrow{\;g\;} Z$
\[ \varphi(g)\, \varphi(f) = \varphi(gf) \]

Le premier membre est
\[
p_{31*}\bigl( p_{32}^{*}(\varphi(g))\, p_{21}^{*}(\varphi(f)) \bigr)
\]
avec
\[
p_{32}^{*}(\varphi(g)) = p_{32}^{*}\bigl( \Gamma_{g*}(1) \bigr) = \lambda_{*}\bigl( \mu^{*}(1) \bigr) = \lambda_{*}(1),
\qquad
p_{21}^{*}(\varphi(f)) = p_{21}^{*}\, \Gamma_{f*}(1) = \lambda'_{*}\bigl( \mu'^{*}(1) \bigr) = \lambda'_{*}(1)
\]
\note{Devant $\lambda_{*}(\mu^{*}(1))$, un \struck{$p_{*}$} est biffé.}
\[
\varphi(g)\, \varphi(f) = p_{31*}\bigl( \lambda_{*}(1)\, \lambda'_{*}(1) \bigr)
\]
Mais
\[
\lambda_{*}(1)\, \lambda'_{*}(1) = \lambda_{*}\bigl( \lambda^{*}( \lambda'_{*}(1) ) \bigr) \quad (\text{form.\ proj.})
\]
\[
\lambda^{*} \lambda'_{*}(1) = \Gamma_{f*}\bigl( \Gamma_{gf}^{*}(1) \bigr) = \Gamma_{f*}(1)
\]
\[
= \lambda_{*}\, \Gamma_{f*}(1) = \nu_{*}(1)
\]
où $\nu \colon X \to X \times Y \times Z$ est donné par
$\nu(x) = (x, f(x), gf(x))$.

\begin{tikzcd}[column sep=small, row sep=small, nodes={font=\scriptsize}]
 & & X \arrow[dl, hook, "\Gamma_{gf}"'] \arrow[dr, hook, "\Gamma_{f}"] & & \\
 & X \times Z \arrow[dl, "\mu' = p_{X}"'] \arrow[dr, hook, "\lambda' = \Gamma_{f} \times Z"] & & X \times Y \arrow[dl, hook, "\lambda = \mathrm{id}_{X} \times \Gamma_{g}"'] \arrow[dr, "\mu = p_{Y}"] & \\
X \arrow[dr, hook, "\Gamma_{f}"'] & & X \times Y \times Z \arrow[dl, "p_{21}"'] \arrow[dr, "p_{32}"] \arrow[dd, "p_{31}"] & & Y \arrow[dl, hook, "\Gamma_{g}"] \\
 & X \times Y & & Y \times Z & \\
 & & X \times Z & &
\end{tikzcd}

\note{Le carré supérieur est marqué « (\uncertain{cartésien}) ». Un premier
croquis du même diagramme, en haut à gauche de la page, n'est pas redessiné.
Sous le diagramme, les images des points :
$x \mapsto (x, gf(x)) \leftarrow (x, f(x), gf(x))$,
$x \mapsto (x, f(x)) \mapsto (x, f(x), gf(x))$, et en marge
\struck{\ill{}} $(x, f(x), z)$, $(x, y, g(y))$, $(x, f(x), gf(x))$.}

Donc
\[
\varphi(g)\, \varphi(f) = p_{31*}\bigl( \nu_{*}(1) \bigr) = \Gamma_{gf*}(1) = \varphi(gf) \qquad \text{OK.}
\]

Définissons le foncteur $\psi$, [i.e.\ \uncertain{posant}
$\psi(X) = A(X)$,] pour $u \in \operatorname{Hom}_{\mathcal{V}^{A}}(X,Y) = A(X \times Y)$,
et $x \in A(X)$, définissons $u(x) \in A(Y)$ par la formule
\struck{$u(x)$}
\[ u(x) = p_{2*}\bigl( u\, p_{1}^{*}(x) \bigr) \]
Prouvons que l'on a bien une \emph{\uncertain{loi} fonctorielle}.

\page{12}

a) Si $X = Y$, prouvons que $\Delta_{X}(x) = x$. \uncertain{Or}

\begin{tikzcd}
 & X \arrow[d, "\delta"'] & \\
 & X \times X \arrow[dl, "p_{1}"'] \arrow[dr, "p_{2}"] & \\
X & & X
\end{tikzcd}

\[
\Delta_{X}(x) = p_{2*}\bigl( \Delta_{X}\, p_{1}^{*}(x) \bigr) = p_{2*}\bigl( \delta_{*}(1)\, p_{1}^{*}(x) \bigr)
\]
--- form.\ de proj.
\[
\delta_{*}(1)\, p_{1}^{*}(x) = \delta_{*}\bigl( 1 \cdot \delta^{*} p_{1}^{*}(x) \bigr),
\qquad \delta^{*} p_{1}^{*}(x) = x
\]
\[
= p_{2*}\, \delta_{*}(x) = x \qquad \text{OK.}
\]

b) Prouvons $v(u(x)) = (v \circ u)(x)$.
\[ X \xrightarrow{\;u\;} Y \xrightarrow{\;v\;} Z \]

\begin{tikzcd}[column sep=small]
 & & X \times Y \times Z \arrow[dl, "\lambda = p_{XY}"'] \arrow[dr, "\mu = p_{YZ}"] & & \\
 & X \times Y \arrow[dl, "p_{1}"'] \arrow[dr, "p_{2}"] & & Y \times Z \arrow[dl, "q_{1}"'] \arrow[dr, "q_{2}"] & \\
X & & Y & & Z
\end{tikzcd}

\note{Devant $q_{2*}$, un \struck{$p_{1}^{*}$} est biffé.}
\[
v(u(x)) = q_{2*}\bigl( v\, q_{1}^{*}(u(x)) \bigr)
= q_{2*}\Bigl( v\, q_{1}^{*}\bigl( p_{2*}( u\, p_{1}^{*}(x) ) \bigr) \Bigr)
\]
[échange]
\[
q_{1}^{*}\, p_{2*}\bigl( u\, p_{1}^{*}(x) \bigr) = \mu_{*}\bigl( \lambda^{*}( u\, p_{1}^{*}(x) ) \bigr) = \mu_{*}\bigl( \lambda^{*}(u)\, p_{X}^{*}(x) \bigr)
\]
\[
= q_{2*}\Bigl( \mu_{*}\bigl( \mu^{*}(v)\, \lambda^{*}(u)\, p_{X}^{*}(x) \bigr) \Bigr)
\]
(proj., attention signe)
\[
= p_{Z*}\bigl( p_{YZ}^{*}(v)\, p_{XY}^{*}(u)\, p_{X}^{*}(x) \bigr)
\]

\note{Le signe d'égalité qui introduit la ligne de l'échange est surmonté d'une
marque \ill{}.}

On a d'autre part

\begin{tikzcd}[column sep=small]
 & X \times Y \times Z \arrow[dl, "p_{XY}"'] \arrow[d, "p_{YZ}"] \arrow[dr, "p_{XZ}"] & & \\
X \times Y & Y \times Z & X \times Z \arrow[dl, "r_{1}"'] \arrow[dr, "r_{2}"] & \\
 & X & & Z
\end{tikzcd}

\[
(v \circ u)(x) = r_{2*}\bigl( (v \circ u)\, r_{1}^{*}(x) \bigr)
= r_{2*}\Bigl( p_{XZ*}\bigl( p_{YZ}^{*}(v)\, p_{XY}^{*}(u) \bigr)\, r_{1}^{*}(x) \Bigr)
\]
\[
p_{XZ*}\bigl( p_{YZ}^{*}(v)\, p_{XY}^{*}(u)\, p_{XZ}^{*}\, r_{1}^{*}(x) \bigr),
\qquad p_{XZ}^{*}\, r_{1}^{*}(x) = p_{X}^{*}(x)
\]
\[
= p_{Z*}\bigl( p_{YZ}^{*}(v)\, p_{XY}^{*}(u)\, p_{X}^{*}(x) \bigr)
\]
OK.

\page{14}

\subsection*{Commutativité du carré de foncteurs}

Il suffit de le voir sur les flèches, donc : prouvons pour
$f \colon X \to Y$ dans $\mathcal{V}$, et $x \in A(X)$
\[ f_{*}(x) = \varphi(f)(x) \]
Or on a

\begin{tikzcd}
 & X \arrow[d, "\Gamma_{f}"] & \\
 & X \times Y \arrow[dl, "p_{1}"'] \arrow[dr, "p_{2}"] & \\
X & & Y
\end{tikzcd}

\[
\varphi(f)(x) = p_{2*}\bigl( \varphi(f)\, p_{1}^{*}(x) \bigr) = p_{2*}\bigl( \Gamma_{f*}(1)\, p_{1}^{*}(x) \bigr)
\]
\[
\Gamma_{f*}(1)\, p_{1}^{*}(x) = \Gamma_{f*}\bigl( \Gamma_{f}^{*}\, p_{1}^{*}(x) \bigr),
\qquad \Gamma_{f}^{*}\, p_{1}^{*}(x) = x
\]
\[
= p_{2*}\, \Gamma_{f*}(x) = x \qquad \text{OK.}
\]

\note{Le dernier membre est écrit $x$ ; le calcul donne $f_{*}(x)$, puisque
$p_{2}\Gamma_{f} = f$.}

\emph{Remarque} \struck{Question} Comment interpréter le foncteur
contravariant $\mathcal{V}^{\circ} \to k$-mod gr.\ défini par $A$,
\uncertain{i.e.}\ composé de $A$ et le foncteur
$\omega^{*} \colon C^{\circ} \to k$-mod gr.\ (associant à $(f_{*}, f^{*})$
l'application $f^{*}$) ? En d'autres termes, comment interpréter
$A(f)^{*} \colon A(Y) \to A(X)$, si $f \colon Y \to X$ ? On démontre que c'est
défini par
\[
\varphi(f)^{*} \in \operatorname{Hom}_{\mathcal{V}^{A}}(Y,X) = A(Y \times X),
\]
déduit de $\varphi(f) \in A(X \times Y)$ par l'isom.\ can.\
\uncertain{voir dessous}. Vérification comme dessus.

\note{« $f \colon Y \to X$ » est bien ce qu'on lit ; pour
$A(f)^{*} \colon A(Y) \to A(X)$ on attend $f \colon X \to Y$.}

On définit un \uncertain{iso.}\ foncteur canonique
\[ s \colon \mathcal{V}^{A\circ} \to \mathcal{V}^{A\circ} \]
induisant l'identité sur les objets, par
\[
u \mapsto {}^{t}u \colon \operatorname{Hom}_{\mathcal{V}^{A}}(X,Y) = A(X \times Y) \to \operatorname{Hom}_{\mathcal{V}^{A}}(Y,X) = A(Y \times X)
\]

\marginal{vérifier \uncertain{que c'est} multiplicatif !}

\note{Les deux membres de $s$ portent l'exposant $A\circ$ sur la page ;
l'usage qui en est fait page 16, $\varphi^{*} = s\varphi^{\circ} \colon
\mathcal{V}^{\circ} \to \mathcal{V}^{A}$, demande un but sans $\circ$. La
phrase se poursuit en tête de la page 16.}

\page{16}

[provenant de] l'isom.\ canonique entre $X \times Y$ et $Y \times X$. Soit
alors
\[ \varphi^{*} = s \varphi^{\circ} \colon \mathcal{V}^{\circ} \to \mathcal{V}^{A} \]
on constate alors que le composé
\[
\psi \varphi^{*} = \psi(s \varphi^{\circ}) = (\psi s) \varphi^{\circ} \colon \mathcal{V}^{\circ} \to k\text{-mod gr.}
\]
est \struck{comme} \emph{\uncertain{égal}} :
$\omega^{*} A^{\circ} \colon \mathcal{V}^{\circ} \to k$-mod gr.

\marginal{\uncertain{à vérifier}}

\struck{Ainsi, $\mathcal{V}^{A}$ apparaît comme une \ill{}
d'\uncertain{enveloppe}}

\emph{Cas où $\mathcal{V}$ a un objet final $e$, et de plus} $A(e) = k$,
[Sinon, on remplace $k$ par $A(e)$ !] Alors les $A(X)$ sont naturellement
\uncertain{jaugés} ; si c'est des alg.\ de Poincaré, il suffit alors de
connaître \struck{ces} \uncertain{ces} alg.\ \uncertain{jaugées}, + les $f^{*}$
($f \colon X \to Y$ dans $\mathcal{V}$) pour connaître $A$.

\note{Le mot lu « jaugés » revient page 18 ; il pourrait aussi être une autre
forme en \emph{-gés} ou \emph{-gués}. Les formes $k_{A(X)}$ qui suivent sont
celles qu'il a en vue.}

Soit
\[ u \in \operatorname{Hom}_{\mathcal{V}^{A}}(X,Y) = A(X \times Y) \]
\[ su \in \operatorname{Hom}_{\mathcal{V}^{A}}(Y,X) = A(Y \times X) \]
d'où
\[ \tilde{u} \colon A(X) \to A(Y) \qquad \widetilde{su} \colon A(Y) \to A(X) . \]

Montrons que [si $u$ \uncertain{central} dans $A(X \times Y)$, \ill{}]
$\tilde{u}$ et $\widetilde{su}$ sont transposés l'un de l'autre, pour les
formes $k_{A(X)}$ et $k_{A(Y)}$ sur $A(X)$, $A(Y)$, i.e.\ on a pour
$x \in A(X)$, $y \in A(Y)$, $u \in A(X \times Y)$
\[
k_{Y}\bigl( y\, u_{*}(x) \bigr) = k_{X}\bigl( u^{*}(y)\, x \bigr)
\]
i.e.
\[
g_{*}\bigl( y\, u_{*}(x) \bigr) = f_{*}\bigl( u^{*}(y)\, x \bigr) .
\]
(Car $u$ central !!) En effet

\begin{tikzcd}
A(X) \arrow[r, bend left=20, "u_{*}"] \arrow[d, "k_{X}"'] & A(Y) \arrow[l, bend left=20, "u^{*}"] \arrow[d, "k_{Y}"] \\
k & k
\end{tikzcd}

\begin{tikzcd}
 & X \times Y \arrow[dl, "p_{1}"'] \arrow[dr, "p_{2}"] & \\
X \arrow[dr, no head, "f"'] & & Y \arrow[dl, no head, "g"] \\
 & e &
\end{tikzcd}

\note{Entre les deux $k$ du premier diagramme, un signe $=$.}

\[
g_{*}\bigl( y\, u_{*}(x) \bigr) = g_{*}\bigl( y\, p_{2*}( u\, p_{1}^{*}(x) ) \bigr) = g_{*}\, p_{2*}\bigl( p_{2}^{*}(y)\, u\, p_{1}^{*}(x) \bigr) = k\bigl( p_{2}^{*}(y)\, u\, p_{1}^{*}(x) \bigr)
\]
\[
f_{*}\bigl( u^{*}(y)\, x \bigr) = f_{*}\bigl( p_{1*}( u\, p_{2}^{*}(y) )\, x \bigr) = f_{*}\, p_{1*}\bigl( u\, p_{2}^{*}(y)\, p_{1}^{*}(x) \bigr) = k\bigl( u\, p_{2}^{*}(y)\, p_{1}^{*}(x) \bigr)
\]
\[
= k\bigl( p_{2}^{*}(y)\, u\, p_{1}^{*}(x) \bigr) \qquad (\text{mod.\ signe !})
\]

\page{18}

\emph{Remarque} On n'a jamais, dans les calculs précédents, utilisé la
structure additive des $A(X)$, seulement leur structure de monoïdes. Donc il y
a lieu de formuler la théorie des correspondances en prenant \struck{des
$X \mapsto$} $A(X)$ \struck{qui sont des} [foncteurs dans les] « monoïdes
bivariants » .... \struck{Il faudrait être \struck{prudent}
\uncertain{rigoureux} pour les signes, dans le cas de monoïdes non
commutatifs, en prenant des monoïdes gradués [les \uncertain{diagrammes}
\ill{}] \ill{} anticommutatifs} Ce point de vue cependant me semble pas très
commode dans le cas où les 2 formules de projection (gauche et droite) ne sont
pas valables l'une et l'autre, et qu'il faut introduire des signes. Il
faudrait plutôt remplacer les monoïdes par des « monoïdes
\uncertain{jaugés} » $(A^{i})_{i \in \mathbf{Z}}$, loi de composition
$A^{i} \times A^{j} \to A^{i+j}$ (NB : L'ensemble somme des $A^{i}$
\struck{est} muni d'une structure de monoïde) et \struck{\ill{}}
\uncertain{applications} \struck{\ill{}} \uncertain{bijectives}
$\sigma^{i} \colon A^{i} \to A^{i}$ satisfaisant $\sigma^{i2} = \mathrm{id}$,
$\sigma^{i} = \mathrm{id}$ si $i$ pair, \struck{pour} \uncertain{enfin chaque}
$A^{*}$ muni d'un \emph{degré} $d(A^{*})$, qui permet d'imposer le degré des
$f_{*}$ dans $(f_{*}, f^{*}) \colon A^{*} \to B^{*}$ comme
$d(B^{*}) - d(A^{*})$. On peut alors imaginer comment les \emph{signes} doivent
être contrôlés dans les formules de projection et d'échange
(\ill{} on \uncertain{sait} que \ill{} sont \ill{}).

\note{Le passage biffé est cancellé en bloc : un trait sur « Il faudrait être
prudent », corrigé en interligne, un encadré autour de la suite et un long trait
ondulé par-dessus ; on n'en donne que ce qui se lit. Au milieu, « (gauche et
droite) » est ajouté en interligne.}

\page{20}

Pour tout $X \in \mathcal{V}$, soit $A'(X) \subset A(X)$, [avec
$1_{A(X)} \in A'(X)$, stables par \uncertain{somme}.] Conditions équiv.

1°) Les $A'(X)$ stables par images inverses, images directes par tout
morphisme de $\mathcal{V}$, et par produit [ou encore, par produits
$\boxtimes$].

2°) \struck{Soit} a) $u \in A'(X \times Y)$, $x \in A'(X)$, alors
$u(x) \in A'(Y)$.

b) Si $x \in A'(X)$, $y \in A'(Y)$, alors $x \boxtimes y \in A'(X \times Y)$

c) Pour tout morphisme $X \to Y$ [il suffit de \ill{}] le graphe
$\gamma_{f} = \Gamma_{f*}(1) \in A(X \times Y)$ est dans $A'(X \times Y)$.

\note{Sous la ligne de c), deux mots, « \ill{} les », ne se rattachent à rien
de sûr. « tout » est entouré dans c) et relié au mot « tout » de 1°) et à la
note marginale.}

\marginal{il faudrait donner un énoncé d'équivalence plus général, où on
remplace \uncertain{« tout »} par \ill{} $H$ morphismes de $\mathcal{V}$ \ill{}}

\struck{NB Les $A'$ correspondent aux sous-catégories (pleines) de
$\mathcal{V}^{A}$ ayant mêmes objets que $\mathcal{V}^{A}$}

\emph{NB} Si $\mathcal{V}$ a un objet final, les $A'$ correspondent exactement
aux sous-catégories de $\mathcal{V}^{A}$ qui ont mêmes objets que
$\mathcal{V}^{A}$, stables par $\boxtimes$ de flèches, et contenant les
\struck{\ill{}} morphismes de $\mathcal{V}$ et leurs $s$ \struck{\ill{}}
${}^{s}\Gamma_{f}$ (il suffit ${}^{s}\delta_{X}$).

\note{Ce second NB est récrit dans un encadré sur le premier ; « mêmes objets
que » est suivi sur la page d'un $A$ seul.}

\emph{Dém} 1°) $\Rightarrow$ 2°)

1° forme simple [avec $\boxtimes$] implique forme non simple, car
\[
x \otimes y = \delta^{*}(x \boxtimes y) \qquad [x, y \in A(X),\ \delta = \mathrm{diag}_{X} \colon X \to X \times X]
\]
\marginal{images inverses par $\delta$}

1° forme non simple implique forme simple, car
\[ x \boxtimes y = p_{1}^{*}(x)\, p_{2}^{*}(y) . \]
\marginal{images inverses par pr.}

a) car
\[ u(x) = p_{2*}\bigl( u\, p_{1}^{*}(x) \bigr) . \]
[\ill{} $=$ pour $u^{*}(y) \in A'(X)$ si $y \in A'(Y)$]
\marginal{images directes, inverses par pr \uncertain{de 2}}

b) clair

c) clair, car $1_{X} \in A'(X)$ par hyp.
\marginal{images directes par graphes}

2°) $\Rightarrow$ 1°). Stabilité par images \struck{\ill{}} directes : car
\[ f_{*}(x) = \gamma_{f}(x) , \]
\note{Un premier $\gamma_{f}$ est biffé : \struck{$\gamma_{f}$}.}
on applique a) et c).

\marginal{produits}

Stabilité par images \struck{directes} inverses : car
$f^{*}(y) = ({}^{t}\gamma_{f})_{*}(y)$, et comme
$\gamma_{f} \in A'(X \times Y)$, $\Rightarrow$
${}^{t}\gamma_{f} \in A'(Y \times X)$ (stabilité par isom.),

Stabilité par $\boxtimes$ claire.

\note{Les notes marginales de cette démonstration sont des étiquettes encadrées
qui disent quelle stabilité sert à quelle étape. La page se termine ici, la
moitié inférieure blanche ; les sous-foncteurs $A'$ sont repris page 22, dans
le lot suivant.}

\end{document}
